   \chapter{ Management de projet }
 \label{chap3}
 \minitoc
 \section{Introduction}
 \label{sec21}

Ce chapitre détaille les éléments de conduite de projet présentés de façon synthétique 
dans le rapport d'action : découpage et suivi des tâches, planning prévisionnel et réalisé, 
démarche et outils, organisation de l'équipe, et missions menées en parallèle du projet 
IPANEMA.
 
\section{Découpage et suivi du projet}

Le projet a été découpé, dès son lancement, en cinq grandes phases elles-mêmes 
déclinées en tâches plus fines, structurant l'ensemble du déroulement jusqu'à la mise en 
production : cadrage et correction du modèle de données, refonte du flux Alteryx, modélisation 
Power BI, développement des mesures DAX et réalisation du dashboard, puis recette et mise 
en production. Ce découpage a été formalisé avec mon maître d'apprentissage sous forme de 
cartes mentales, et suivi au fil du projet à l'aide d'un tableau Kanban, dont une vue simplifiée 
est présentée ci-dessous.

\begin{figure}
    \centering
    \roundedfullwidthimage{cover/figures/Vue_simplifiée_du_tableau_Kanban_de_suivi.png}
    \caption{Vue simplifiée du tableau Kanban de suivi du projet IPANEMA.}
    \label{fig:Vue_simplifiée_du_tableau_Kanban_de_suivi}
\end{figure}

À la date de rédaction de ce mémoire, seules les tâches liées à la recette et à la mise en 
production restent au statut « À faire », ce qui traduit un projet en phase finale, dont le 
développement est intégralement terminé.

\section{Planning prévisionnel et réalisé — analyse des écarts}

Dès le lancement du projet, le 23 février 2026, j'ai construit avec Laurent un planning 
prévisionnel détaillé pour les quatre phases de réalisation décrite dans 1., avec les échéances 
suivantes : fin du cadrage et de la correction du modèle de données le 13 mars, fin de la 
refonte du flux Alteryx le 24 avril, fin de la modélisation Power BI et du développement des 
mesures DAX le 5 juin, puis fin de la réalisation du dashboard et des validations progressives 
le 13 juillet. Ce travail de planification en amont, mené conjointement avec mon maître 
d'apprentissage, a permis d'anticiper la charge de chaque phase, de sécuriser les jalons du 
projet et de fiabiliser le suivi d'avancement tout au long des cinq mois de développement. \\

Dans les faits, le planning réalisé a suivi ce découpage jusqu'au 13 juillet, avant de 
s'écarter du prévisionnel sur un point notable : le projet a été mis en stand-by du 13 au 29 
juillet 2026, période durant laquelle j'ai été mobilisée sur une mission au Brazilian Innovation 
Center (BIC) de Pierre Fabre à Rio de Janeiro. Cette interruption a mécaniquement décalé la 
reprise des travaux et, combinée aux contraintes de disponibilité des parties prenantes en 
période estivale, a conduit à reporter l'ouverture des ateliers de recette au-delà de la date 
initialement prévue.

\begin{figure}[H]
    \centering
    \roundedfullwidthimage{cover/figures/Planning_prévisionnel_vs_réalisé.png}
    \caption{Planning prévisionnel vs réalisé, avec l'écart lié à la mission BIC.}
    \label{fig:Planning_prévisionnel_et_réalisé}
\end{figure}

Cet écart illustre une contrainte propre au contexte d'une alternance : le planning d'un 
projet doit composer avec d'autres engagements de l'apprenti au sein de l'entreprise. Il n'a pas 
remis en cause la faisabilité du projet, mais a nécessité un ajustement du calendrier de recette, 
discuté et validé avec Laurent et l'équipe Qualité à la reprise.

\section{Démarche et outils}

Le projet a été conduit selon une démarche agile adaptée au fonctionnement de l'équipe, 
plutôt qu'un cycle en V strict : plutôt que de figer l'ensemble des spécifications avant tout 
développement, le projet a démarré par la correction du modèle de données en lien avec le 
prestataire, puis a progressé par itérations successives, structurées par la boucle 

\begin{center}
\small\sffamily
\textbf{Planifier} $\longrightarrow$ 
\textbf{Réaliser} $\longrightarrow$ 
\textbf{Tester} $\longrightarrow$ 
\textbf{Recueillir les retours} $\longrightarrow$ 
\textbf{Ajuster}
\end{center}

Cette boucle s'est concrétisée par des tests et revues réalisés progressivement pendant le 
développement avec l'interlocuteur métier, plutôt qu'une validation reportée à la seule fin du 
projet. \\

Le suivi opérationnel s'est appuyé sur le tableau Kanban présenté au §1, complété par 
des cartes mentales pour la structuration des tâches. Des réunions bimensuelles réunissant 
l'équipe métier (équipe Qualité) et Laurent ont rythmé le projet, permettant de partager 
l'avancement, valider les orientations et définir les prochaines étapes. Entre ces points, j'ai 
conduit mes travaux en autonomie ; en cas de blocage, je sollicitais directement Laurent, qui 
m'aidait à le résoudre ou m'orientait vers la personne compétente lorsque je ne l'avais pas 
encore identifiée moi-même. \\

Sur le plan technique, le développement s'est appuyé sur Power BI Desktop et Alteryx 
Designer en environnement local, avec publication respective sur Power BI Service et Alteryx 
Server pour l'exécution en production. Les données sources (listes SharePoint générées par 
l'application IPANEMA) et les échanges d'équipe se sont appuyés sur l'écosystème Microsoft 
365 (SharePoint, Teams). 

\section{Organisation et relations avec l'équipe}

Le projet s'est déroulé dans un contexte de réorganisation de la R\&D DCPC. Avant cette 
réorganisation, l'équipe IT-Data R\&D DCPC formait un groupe unique, de taille plus réduite, 
rattaché directement à Yves Labourdique. Depuis, cette équipe a été redéployée au sein d'une 
nouvelle Direction Portfolio et Transformation Digitale, structurée en trois services placés 
chacun sous la responsabilité d'un manager dédié : Transformation / Data et Portefeuille, 
Digital Product Science, et Digital Product Dev et Perf — mon équipe, dirigée par Audrey 
Olivan et composée d'une dizaine de collaborateurs.

\begin{figure}[H]
    \centering
    \roundedfullwidthimage{cover/figures/Organigramme_avant_après_la_réorganisation.png}
    \caption{Organigramme avant / après la réorganisation.}
    \label{fig:Organigramme_avant_après_la_réorganisation}
\end{figure}


Cette réorganisation répondait à l'ambition de regrouper les compétences de gestion de 
portefeuille, de digital, de data et d'intelligence artificielle au sein d'une direction unique, afin 
de renforcer l'accompagnement des services de la R\&D, de mutualiser les compétences et de 
favoriser le développement des expertises. Pour le projet IPANEMA, cette transition n'a pas 
modifié le cadrage opérationnel assuré par Laurent, mais elle a fait évoluer le contexte 
managérial dans lequel le projet s'inscrivait, avec un rattachement désormais clarifié au sein 
d'un pôle Digital Product Dev et Perf explicitement dédié aux produits et outils digitaux. \\

Au-delà de la relation avec mon maître d'apprentissage, le projet a nécessité une 
collaboration directe et régulière avec l'équipe Qualité, commanditaire et utilisatrice finale du 
dashboard : réunion de recueil du besoin en amont, tests et revues progressifs pendant le 
développement, échange sur la liste des personnes autorisées à accéder au dashboard, et 
définition conjointe des scénarios de recette. \\

Cette dynamique d'équipe s'est notamment illustrée lors d'une journée de département 
réunissant l'ensemble des collaborateurs autour d'activités de team building et d'un temps 
d'échange collectif sur les objectifs de l'année à venir et les axes stratégiques du groupe.

\begin{figure}[H]
    \centering
    \roundedfullwidthimage{cover/figures/equipe1.jpg}
    \caption{Journée de département, 10 juillet 2026.}
    \label{fig:Journée_de_département}
\end{figure}


\section{Missions parallèles}

En complément du projet IPANEMA, plusieurs missions ont été menées en parallèle au 
cours de cette première année d'alternance, illustrant la diversité des contextes rencontrés. \\

Une mission de trois semaines a été réalisée au Brazilian Innovation Center (BIC) de 
Pierre Fabre à Rio de Janeiro, du 13 au 29 juillet 2026 — mission à l'origine de l'écart de 
planning décrit au §2. Elle a couvert un inventaire complet des dashboards et bases de 
données du BIC, le déploiement d'un dashboard Power BI de gestion des fournisseurs (traduit 
du portugais vers l'anglais), ainsi qu'une présentation d'acculturation à l'intelligence artificielle 
pour l'équipe locale.\\

En parallèle, j'ai contribué au flux de sourcing destiné à être intégré au système Raw 
Material Insights (présenté en détail dans mon rapport du premier semestre), en développant 
un pipeline Python intégré à Alteryx permettant de consolider des formulaires Excel multi
onglets issus des fournisseurs. \\

J'ai également corrigé un flux cassé pour l'équipe analytique, et développé un dashboard 
pour l'équipe caractérisation. \\

Enfin, j'ai coécrit avec Laurent les spécifications du projet Préfo Actifs[3], une étude de 
faisabilité (POC) portée par le service d'innovation galénique. Ce service évalue l'introduction 
de nouveaux principes actifs dans les formules galéniques au moyen d'analyses de risques, 
historiquement ralenties par la dispersion des données nécessaires entre plusieurs systèmes 
d'information et fichiers de travail. L'objectif du POC est de modéliser ces données de 
formulation en vue d'analyses prédictives et prescriptives, avec la possibilité, à terme, d'un 
agent conversationnel capable de répondre à des questions métier ciblées (compatibilité entre 
actifs, mode d'introduction, plage de pH...). \\

Ma contribution a porté sur l'identification et la cartographie des sources de données 
mobilisables : la table i-formula de la Data Market Place (DMP), les données de formule issues 
de Teexma, ainsi que plusieurs fichiers de travail SharePoint (dossiers actifs, analyse de 
risques Innothèque - ADRI, et retours d'expérience - REX). Cette cartographie a permis de 
préparer la construction d'un flux Alteryx consolidant ces sources hétérogènes vers un dossier 
de sortie documenté, en cohérence avec le dictionnaire de données du projet. 

\begin{figure}[H]
  \centering
  \roundedfullwidthimage{cover/figures/Sources_de_données_identifiées_pour_le_projet_Préfo_Actifs.png}
  \caption{Sources de données identifiées pour le projet Préfo Actifs.}
  \label{fig:Sources_de_données_identifiées_pour_le_projet_Préfo_Actifs}
\end{figure}

Le projet est organisé en lots pour permettre une livraison progressive : une première 
vue synthétique, puis l'ensemble des vues avec le template graphique du groupe, avant une 
amélioration itérative de l'ergonomie et des fonctionnalités — une démarche de priorisation 
similaire à celle retenue pour IPANEMA. \\

Prises ensemble, ces missions parallèles montrent une même logique transposée à des 
contextes différents : recueillir un besoin métier, cartographier des sources de données 
hétérogènes, et construire une chaîne de traitement fiable jusqu'à la restitution. Elles ont 
également contribué à consolider ma maîtrise d'Alteryx et de Power BI au-delà du seul 
périmètre d'IPANEMA, et à diversifier mes interlocuteurs au sein de la R\&D DCPC (équipes 
Sourcing, Analytique, Caractérisation, Innovation galénique), en complément de l'équipe 
Qualité.

\section{Conclusion}

Au total, ce chapitre a montré une conduite de projet structurée dès le lancement, un 
planning tenu à un aléa près assumé et anticipé dans son traitement, une démarche itérative 
outillée par le Kanban et les cartes mentales, et une équipe dont l'organisation a évolué en 
cours de projet sans remettre en cause le cadrage opérationnel du projet IPANEMA. Le 
chapitre suivant revient sur le bilan personnel de cette première année d'alternance.
